% !TeX root = ../../main.tex
% ------------------------------------------------------------------
%  4.2.1 Especificaciones de los implementos a proveer (hardware y software).
%
% ------------------------------------------------------------------

\section{Especificaciones Implementos a proveer (Hardware y Software)}
\label{sec:implementos}

Esta sección resume y analiza el equipamiento y el software que la solución requiere en las instalaciones del CLIENTE, en el terreno y en la plataforma de nube. El detalle de cada elemento, con su producto ofertado, su ubicación, su cantidad y su justificación, se entrega en el Formulario T-11. El hardware de terreno ---terminales, lectores, impresoras portátiles y sensores--- lo adquiere el CLIENTE, y el adjudicatario lo especifica (PUCV, 2026a, cap. 11, p. 20). La infraestructura on-premise la provee el adjudicatario dentro del contrato (PUCV, 2026c, art. 14.2, p. 10). Todo el equipamiento lo instala, integra y mantiene el adjudicatario.

\subsection{Síntesis del equipamiento por familia}

La Tabla~\ref{tab:t26} cruza cada familia de equipamiento con el sitio donde se instala. Las cifras son las unidades en operación; la reserva se analiza a continuación de la tabla.

\begin{tablalafrox}{Equipamiento en operación por familia y por sitio}{tab:t26}%
  {Y H{0.13} H{0.14} H{0.17} H{0.16}}%
  {\cab{Familia} & \cab{Talca} & \cab{Concepción} & \cab{Cada cross-docking} & \cab{Terreno}}
  Servidores y nodos de cómputo & 3 & 1 & 1 & -- \\
  Respaldo local (NAS WORM) & 1 & -- & -- & -- \\
  Firewalls & 2 & 2 & 2 & -- \\
  Switches de núcleo, de gestión y de sitio & 3 & 2 & 2 & -- \\
  Switches de acceso en las bodegas & 4 & 2 & -- & -- \\
  UPS de gabinetes de piso & 4 & 2 & -- & -- \\
  UPS de gabinetes de borde & -- & 1 & 1 & -- \\
  Enlaces satelitales D-06 & 1 & 1 & 1 & -- \\
  Puntos de acceso Wi-Fi 6E & 39 & 20 & 2 & -- \\
  Terminales de bodega & 120 & 60 & 2 & -- \\
  Impresoras de andén & 4 & 2 & -- & -- \\
  Balanzas de recepción & 2 & 1 & -- & -- \\
  Puntos de temperatura en cámaras & 15 & 6 & -- & -- \\
  Gateways IoT & 2 & 1 & -- & -- \\
  Estaciones de trabajo & 5 & 3 & -- & -- \\
  Terminales de preventa & -- & -- & -- & 62 \\
  Equipos de reparto por camión & -- & -- & -- & 96 \\
  Termógrafos de camión & -- & -- & -- & 28 \\
\end{tablalafrox}

{\footnotesize Fuente: elaboración propia.\par}

La tabla incluye cinco enlaces satelitales: uno en Talca, uno en Concepción y uno por cada cross-docking.

La tabla muestra que la operación se concentra en el terreno y en las bodegas, no en la sala técnica. Solo en la calle trabajan 62 terminales de preventa y 96 camiones equipados, cada uno con un terminal de conductor, una impresora de cabina y un terminal de pago; en las bodegas y en los cross-docking, 186 terminales y los puntos de acceso que los conectan. Frente a eso, el cómputo y el almacenamiento suman 8 equipos: 3 servidores en Talca, 1 en Concepción, 3 mini-PC y la NAS. Esa proporción refleja el caso, donde la operación ocurre en la calle, en la cámara y en el andén, y explica por qué la gestión de dispositivos (N-13) es un componente con emplazamiento propio y no un accesorio.

Los equipos de reparto se cuentan por camión y no por conductor, porque los conductores de los transportistas rotan sin aviso y el equipo queda en el vehículo (S-30). Los terminales de bodega se comparten entre turnos, como contempla el perfil operacional del caso (RT-12.11; PUCV, 2026a, cap. 15, p. 27), por lo que su cantidad la fija el turno con más personas trabajando a la vez: los 120 preparadores nocturnos de Talca y los 60 de Concepción (S-39); la carga de los camiones la hace esa misma cuadrilla (S-33). Los puntos de acceso son una estimación preliminar por superficie, con mayor densidad dentro de las cámaras, que el estudio de cobertura confirma (S-36).

A las unidades en operación se suma la reserva del numeral 8.4 de las Bases Técnicas Transversales (PUCV, 2026b, cap. 8, p. 19): el 10 \% del parque de cada tipo de dispositivo y de componente crítico, redondeado hacia arriba. Con ella se cotizan 69 terminales de preventa, 106 de cada equipo de reparto, 205 terminales de bodega, 72 puntos de acceso, 7 impresoras de andén, 4 balanzas, 7 módulos y 24 sondas de temperatura, 31 termógrafos y 4 gateways IoT, además de una unidad de reserva por cada modelo de equipo de red y para el mini-PC: una para los firewalls de los centros de distribución, una para los de los cross-docking, una para los switches de núcleo, una para los de acceso, una para los seis switches de los cross-docking (6 + 1 = 7) y una para el mini-PC. El switch de gestión no lleva reserva, porque su falla no detiene la operación. Los servidores no llevan una unidad entera de reserva: se cubren con repuestos por componente, como discos y fuentes. A tres años, el crecimiento de los viajes y de la dotación que proyecta el caso agrega 15 unidades de cada equipo de reparto, 8 terminales de preventa y 28 terminales de bodega (S-31, S-32); los termógrafos no crecen, porque el caso no proyecta más camiones con equipo de frío (S-38). El detalle elemento por elemento está en el Formulario T-11, el cálculo de cada cantidad en el Anexo 4-W y los supuestos en el registro del Subdocumento 3.

Se proveen ocho estaciones nuevas para despacho, administración, planificación, calidad y TI. Los equipos existentes de los usuarios de oficina se incorporan a la gestión central con CrowdStrike Falcon, cifrado de disco, parches y control de extraíbles como condición de acceso por Verified Access conforme a ADR-16 (Anexo 4-O).

\subsection{Criterios de selección}

Cada familia se especifica contra una condición del caso que el equipamiento de oficina no resiste:

\begin{itemize}
  \item Los 22 terminales de la cuadrilla de congelado están certificados para trabajar a $-$30\,\textdegree{}C con guantes gruesos y batería de recambio en caliente, porque la preparación de congelados ocurre a $-$22\,\textdegree{}C. Los 183 restantes atienden la bodega seca, la cámara de refrigerado sobre 0\,\textdegree{}C y los cross-docking, y son el modelo estándar, porque esas condiciones no exigen la versión para congelado (S-34).
  \item Los terminales de terreno operan un turno completo sin señal, a una mano y bajo sol directo.
  \item Los mini-PC de los cross-docking son de rango industrial, porque las naves no tienen climatización y el supuesto de diseño es de hasta 50\,\textdegree{}C bajo la techumbre en verano.
  \item Los 28 termógrafos en operación cubren los 18 camiones propios y los 10 refrigerados de los transportistas que declara el caso, porque el instrumento viaja con la carga y su registro ampara la cadena de frío aun cuando el vehículo no sea del CLIENTE; en los camiones de terceros se instala con el acuerdo de cada transportista (S-28). Están calibrados contra un patrón NIST en dos puntos, porque su registro es la prueba de la cadena de frío ante un cliente o ante la autoridad sanitaria.
\end{itemize}

En la sala técnica se aplica redundancia sin sobrecompra:

\begin{itemize}
  \item Talca: los tres nodos idénticos disponen cada uno de 32 hilos, 64 GB RAM y dos NVMe de 960 GB para Ceph sin RAID. Ceph mantiene tres réplicas y quórum con dos nodos; la configuración cubre la carga a 3× con un nodo caído, según el Anexo 4-W.
  \item Concepción: el servidor se dimensiona para su bodega y los switches de núcleo operan en par conforme a RT-08.03 (PUCV, 2026b, cap. 8, p. 18).
  \item Los cinco sitios: los firewalls van en par activo/pasivo; los cross-docking disponen además de dos switches industriales por plataforma. La reposición del mini-PC se apoya en la unidad de reserva conservada en Talca.
\end{itemize}

Los equipos de infraestructura cumplen RT-08.04 (PUCV, 2026b, cap. 8, p. 18) con doble fuente y conexión a circuitos distintos. Los terminales móviles tienen batería, los puntos de acceso reciben PoE de switches con doble fuente, y las impresoras, balanzas y estaciones de trabajo se cubren mediante unidades alternativas y circuitos protegidos.

Todo el equipamiento es nuevo, sin uso previo y con garantía de fábrica vigente desde la recepción conforme (RT-08.06; PUCV, 2026b, cap. 8, p. 18). Su vida útil esperada es de cinco años, mayor que los 56 meses del contrato. Los repuestos provienen de la reserva del 10\,\% y de la garantía del fabricante. Cada unidad dada de baja se repone desde esa reserva (RT-08.13; PUCV, 2026b, cap. 8, p. 19).

\subsection{Software y licenciamiento}

El software de base que opera el adjudicatario es de código abierto y se organiza así:

\begin{itemize}
  \item PostgreSQL con PostGIS, RabbitMQ, Keycloak, Proxmox VE, Ceph con tres réplicas y Docker sostienen la plataforma, junto con el motor de optimización de rutas.
  \item La aplicación está escrita en PHP 8.5 con Laravel 13 y sus dependencias fijadas por Composer. Las dependencias de PHP se revisan en cada construcción con \texttt{composer audit} y su licencia se verifica contra una lista de licencias admitidas antes de incorporarlas.
  \item Los portales usan Angular y TypeScript; las aplicaciones de terreno usan Kotlin.
\end{itemize}

Los únicos elementos con suscripción son los servicios de la plataforma de nube, descritos en la sección~\ref{sec:servicios-nube}, la gestión de dispositivos, los agentes de detección en endpoints y el orquestador de integración continua GitLab CI. Esta composición mantiene la reversibilidad de la solución y evita que el CLIENTE quede atado a un licenciamiento propietario.
